\section{Realización, balance armónico y continuación causal}
\label{part:theory}

\subsection{Fundamentos de la localización}
\label{sec:theory-attribution}

La forma de Lur'e, la función descriptiva, el balance armónico, el criterio de
Matignon, la equivalencia Caputo--Volterra, el algoritmo ABM, la continuación
de bifurcaciones y las realizaciones multivariables poseen literatura previa
\cite{Leonov2013HiddenAttractors,LeonovKuznetsov2013,Matignon1996,
DiethelmFordFreed2002,bonani1999,rogov2019}. La autorreproducción afín de
atractores fraccionarios se estudia en \cite{sayed2020}.

La localización combina las siguientes operaciones:
decidir antes de la búsqueda si una realización escalar es
exacta; cambiar de ruta cuando existe una obstrucción demostrable; extender el
cierre a diccionarios multicanal; preservar una historia causal de Caputo; y
mantener separados álgebra, localización, dinámica y evidencia finita de
cuenca.

\begin{longtable}{L{0.28\linewidth}L{0.27\linewidth}L{0.37\linewidth}}
\caption{Fundamentos y función de cada construcción.}
\label{tab:novelty-map}\\
\toprule
Elemento & Fundamento & Función matemática\\
\midrule
\endfirsthead
\toprule
Elemento & Fundamento & Función matemática\\
\midrule
\endhead
Certificado de rango uno & Geometría de realizaciones Lur'e & Distinguir una
realización escalar exacta de una aproximación.\\
Saltos PWA compartidos & Realimentación escalar PWA & Identificar direcciones
comunes de entrada y salida en las regiones afines.\\
Equivalencia afín & Regla lineal de Caputo y antecedente
\cite{sayed2020} & Transportar el campo y las vecindades de cuenca.\\
Diccionario MIMO & Lur'e y balance multivariable
\cite{bonani1999,rogov2019} & Construir semillas multiharmónicas.\\
Historia única de Caputo & Ecuación de Volterra & Conservar las contribuciones
anteriores durante la continuación.\\
Machado \(N^\mu\) & Función descriptiva fraccionaria publicada
\cite{Machado2015FDF} & Separar residuo auxiliar y físico en cada semilla.\\
Puntos perpetuos enteros & Condición \(J_f f=0\) y localización heurística
\cite{Prasad2015Perpetual,DudkowskiPrasadKapitaniak2017} & Se integran como
una ruta de semillas, con verificación dinámica posterior.\\
FPP--A/FPP--H & Volterra, expansión Picard y operador Caputo secuencial &
Separar estado de arranque y evento con historia.\\
Separación continuación/congelación/reset & No localidad de Caputo &
Separación metodológica entre continuación, congelación del campo y reinicio.\\
Protocolo L0--L6 y niveles EV--TG0--EV--TG7 & Definición de atractor oculto y
validación numérica & Distinguir semilla, recurrencia, cuenca y certificación.\\
\bottomrule
\end{longtable}

\subsection{Certificado de realizabilidad escalar por rango uno}
\label{sec:rank-one-certificate}

Sea \(\Omega\subset\mathbb R^n\) abierto y conexo, y
\(f\in C^1(\Omega,\mathbb R^n)\). Fijados \(A\in\mathbb R^{n\times n}\),
\(b,r\in\mathbb R^n\setminus\{0\}\), escribimos
\begin{equation}
 g(x)=f(x)-Ax,\qquad
 P_b^\perp=I-\frac{bb^{\mathsf T}}{b^{\mathsf T}b},\qquad
 P_r^\perp=I-\frac{rr^{\mathsf T}}{r^{\mathsf T}r}.
 \label{eq:th-projectors}
\end{equation}

\begin{proposition}[necesidad]
\label{prop:rank-one-necessary}
Si existe \(\psi\in C^1(r^{\mathsf T}\Omega,\mathbb R)\) tal que
\[
 f(x)=Ax+b\,\psi(r^{\mathsf T}x)
\]
en \(\Omega\), entonces
\begin{equation}
 P_b^\perp g(x)=0,\qquad
 Dg(x)P_r^\perp=0,\qquad
 \operatorname{rank}Dg(x)\leq1
 \label{eq:th-rank-necessary}
\end{equation}
para todo \(x\in\Omega\).
\end{proposition}

\begin{proof}
La primera igualdad sigue de
\(g(x)=b\psi(r^{\mathsf T}x)\in\operatorname{span}\{b\}\). Por la regla de la
cadena,
\[
 Dg(x)=b\,\psi'(r^{\mathsf T}x)\,r^{\mathsf T}.
\]
Como \(r^{\mathsf T}P_r^\perp=0\), se obtiene
\(Dg(x)P_r^\perp=0\). La matriz exterior \(br^{\mathsf T}\) tiene rango uno,
y multiplicarla por un escalar no puede aumentar el rango.
\end{proof}

La condición de rango por sí sola no basta: una dirección de salida o de
entrada que varíe con \(x\) también puede producir Jacobianos de rango uno. La
prueba debe conservar \(b\) y \(r\) fijos.

\begin{proposition}[suficiencia sobre fibras conexas]
\label{prop:rank-one-sufficient}
Supóngase que cada fibra
\(\Omega_s=\{x\in\Omega:r^{\mathsf T}x=s\}\) es conexa o vacía. Si
\[
 P_b^\perp g(x)=0,\qquad Dg(x)P_r^\perp=0
\]
en \(\Omega\), entonces existe una función
\(\psi:r^{\mathsf T}\Omega\to\mathbb R\) tal que
\[
 f(x)=Ax+b\,\psi(r^{\mathsf T}x).
\]
\end{proposition}

\begin{proof}
La primera igualdad permite definir
\[
 \lambda(x)=\frac{b^{\mathsf T}g(x)}{b^{\mathsf T}b},
 \qquad g(x)=b\lambda(x).
\]
Derivando, \(Dg=b\,D\lambda\). Como \(b\neq0\), la segunda hipótesis implica
\(D\lambda(x)P_r^\perp=0\). Sea \(x(\tau)\) una curva \(C^1\) contenida en una
fibra. Entonces \(r^{\mathsf T}\dot x(\tau)=0\), es decir,
\(\dot x=P_r^\perp\dot x\), y
\[
 \frac{d}{d\tau}\lambda(x(\tau))
 =D\lambda(x(\tau))P_r^\perp\dot x(\tau)=0.
\]
Así, \(\lambda\) es constante en cada fibra conexa. Definir
\(\psi(s)=\lambda(x)\) para cualquier \(x\in\Omega_s\) es, por tanto, no
ambiguo, y \(g(x)=b\psi(r^{\mathsf T}x)\).
\end{proof}

\begin{proposalbox}
\textbf{Criterio de admisión.} Usar
\(P_b^\perp g\) y \(DgP_r^\perp\) para comprobar las identidades de
\cref{prop:rank-one-sufficient} antes de resolver balance armónico. La ruta D0
exacta requiere justificar ambas identidades en todo el dominio y la conexión
de sus fibras. Residuos numéricos pequeños definen una aproximación cuya norma,
dominio y error deben declararse; una tolerancia no convierte esa aproximación
en D0 exacta. Si las identidades fallan, se intenta una transformación T1 o una
realización M1, o se registra explícitamente el modelo escalar aproximado.

\textbf{Prueba de falsación.} Aplicar el certificado a una familia con
Hessiano no lineal de rango mayor que uno. El jerk cuadrático de
\cref{sec:cand-jerk} debe ser rechazado por D0 y aceptado por M1.
\end{proposalbox}

\begin{proposition}[defecto uniforme de una aproximación escalar]
\label{prop:rank-one-approximation-bound}
Sea \(D\subset\Omega\) un dominio de evaluación. Supóngase que
\[
 \sup_{x\in D}\|P_b^\perp g(x)\|_2\le\varepsilon_b,
 \qquad
 \sup_{x\in D}\|Dg(x)P_r^\perp\|_{2\to2}\le\varepsilon_r.
\]
Para cada fibra no vacía \(D_s=\{x\in D:r^{\mathsf T}x=s\}\), elíjase
\(x_s\in D_s\). Supóngase que cada \(x\in D_s\) puede unirse a \(x_s\)
por una curva a trozos \(C^1\) contenida en \(D_s\), de longitud a lo sumo
\(L_{\rm fib}<\infty\), uniforme en \(s\) y \(x\). Entonces
\[
 \psi_D(s)=\frac{b^{\mathsf T}g(x_s)}{b^{\mathsf T}b}
\]
satisface
\[
 \sup_{x\in D}\|f(x)-Ax-b\psi_D(r^{\mathsf T}x)\|_2
 \le\varepsilon_b+L_{\rm fib}\varepsilon_r.
\]
\end{proposition}

\begin{proof}
Escriba \(Q_b=I-P_b^\perp\). Para \(x\in D_s\), el defecto es
\(P_b^\perp g(x)+Q_b[g(x)-g(x_s)]\). A lo largo de una curva \(\gamma\)
contenida en la fibra, \(\gamma'=P_r^\perp\gamma'\); por tanto,
\[
 \|g(x)-g(x_s)\|_2
 \le\int\|Dg(\gamma)P_r^\perp\|_{2\to2}\|\gamma'\|_2
 \le L_{\rm fib}\varepsilon_r.
\]
La cota sigue de \(\|Q_b\|_{2\to2}=1\) y la desigualdad triangular.
\end{proof}

El enunciado controla el defecto del campo en \(D\); no demuestra por sí solo
regularidad de la sección elegida \(s\mapsto x_s\), equivalencia dinámica ni
proximidad de atractores. Si se usa \(\psi_D\) como un nuevo campo, deben
justificarse además su regularidad y el control de sus soluciones. Un muestreo
finito no establece las dos cotas uniformes de la proposición sin un control
adicional entre muestras.

La necesidad del dominio y de \(L_{\rm fib}\) se ve con
\(g(x,y)=(\varepsilon y,0)\), \(A=0\) y \(b=r=e_1\) en
\(\mathbb R^2\). El residual de salida es cero y el transversal vale
\(\varepsilon\), por pequeño que sea; sin embargo, \(g\) no depende sólo de
\(x\), y para toda \(\psi\) el error uniforme respecto de
\(b\psi(x)\) es infinito cuando \(\varepsilon>0\).

\subsection{Criterio de saltos compartidos en campos PWA}
\label{sec:pwa-shared-jumps}

Considérese una partición en regiones \(\Omega_j\) y un campo afín por partes
\[
 f_j(x)=A_jx+d_j.
\]
Elegimos una región base \(j=0\).

\begin{proposition}[realización PWA escalar exacta]
\label{prop:pwa-shared}
Si existen \(b,r\neq0\) y escalares \(\gamma_j,\delta_j\) tales que
\begin{equation}
 A_j-A_0=b\gamma_jr^{\mathsf T},\qquad
 d_j-d_0=b\delta_j
 \label{eq:th-pwa-jumps}
\end{equation}
para toda región, entonces
\begin{equation}
 f_j(x)=A_0x+d_0+b\bigl(\gamma_jr^{\mathsf T}x+\delta_j\bigr).
 \label{eq:th-pwa-realization}
\end{equation}
Si, además, la etiqueta regional \(j\) depende sólo de
\(\sigma=r^{\mathsf T}x\), existe una característica escalar PWA
\(\psi(\sigma)\) que reproduce todo el campo.
\end{proposition}

\begin{proof}
Sustituir \cref{eq:th-pwa-jumps} en
\(A_jx+d_j\) da \cref{eq:th-pwa-realization}. Si cada región es la preimagen
de un intervalo \(I_j\) por \(r^{\mathsf T}x\), se define
\(\psi(\sigma)=\gamma_j\sigma+\delta_j\) para \(\sigma\in I_j\).
\end{proof}

Las condiciones son también necesarias para una realización escalar que use
el mismo \(A_0,d_0,b,r\) y una \(\psi\) afín en cada intervalo. La continuidad
en una frontera \(\sigma=s_j\) añade
\[
 \gamma_{j^-}s_j+\delta_{j^-}
 =\gamma_{j^+}s_j+\delta_{j^+}.
\]
Una discontinuidad controlada puede admitirse, pero entonces debe declararse
la solución diferencial utilizada.

\begin{proposalbox}
\textbf{Aplicación.} Las diferencias entre matrices regionales deben tener
rango uno y compartir direcciones de entrada y salida. La geometría de las
regiones debe depender del mismo observable \(r^{\mathsf T}x\). El diodo
PWL de Chua satisface esta estructura; superficies de conmutación incompatibles
con ese observable impiden la realización escalar.
\end{proposalbox}

\subsection{Equivalencia afín y transporte correcto de cuencas}
\label{sec:affine-caputo}

Sea \(x=Tz+c\), con \(T\) invertible y \(c\) constante. Para Caputo
\({}^{\mathrm C}D^q c=0\) y
\[
 {}^{\mathrm C}D^q x=T\,{}^{\mathrm C}D^q z.
\]
Si
\[
 {}^{\mathrm C}D^q x=Ax+b\psi(r^{\mathsf T}x),
\]
entonces
\begin{equation}
 {}^{\mathrm C}D^q z
 =T^{-1}ATz+T^{-1}(Ac)
 +T^{-1}b\,
 \psi\!\left(r^{\mathsf T}Tz+r^{\mathsf T}c\right).
 \label{eq:th-affine-caputo}
\end{equation}
Para una transformación puramente lineal \(c=0\),
\[
 A_z=T^{-1}AT,\quad b_z=T^{-1}b,\quad r_z=T^{\mathsf T}r.
\]

La transformación no autoriza a conservar esferas numéricamente iguales. Una
bola en coordenadas físicas,
\[
 \|x-x_e\|_2\leq\rho,
\]
se convierte en el elipsoide
\begin{equation}
 (z-z_e)^{\mathsf T}T^{\mathsf T}T(z-z_e)\leq\rho^2.
 \label{eq:th-ellipsoid}
\end{equation}
Usar una bola de radio \(\rho\) en \(z\) cambia el experimento salvo cuando
\(T\) es ortogonal. Esta observación es esencial para comparar cuencas antes y
después de una transformación T1.

\begin{remark}
La literatura incluye transformaciones afines para autorreproducción de
atractores ocultos \cite{sayed2020}. La formulación evaluada transporta además
el protocolo de ocultedad y las tolerancias.
\end{remark}

\subsection{Realización polinómica multicanal y minimalidad}
\label{sec:mimo-dictionary}

Cuando el certificado escalar falla, se usa
\begin{equation}
 {}^{\mathrm C}D^q x=Ax+d+B\Phi(Cx),
 \qquad
 W(z)=C(zI-A)^{-1}B.
 \label{eq:th-mimo}
\end{equation}
\(\Phi=(\phi_1,\ldots,\phi_m)^{\mathsf T}\) es un diccionario declarado de
monomios, funciones PWA o funciones suaves. La verificación exacta es
\[
 R(x)=f(x)-Ax-d-B\Phi(Cx)\equiv0.
\]

Para polinomios, un diccionario completo de monomios reproduce siempre el
campo; el problema científico es encontrar uno pequeño y estructuralmente
útil. Dos identidades permiten reorganizar términos cruzados:
\begin{align}
 uv&=\frac14\bigl[(u+v)^2-(u-v)^2\bigr],
 \label{eq:th-square-polarization}\\
 uvw&=\frac1{24}\bigl[(u+v+w)^3+(u-v-w)^3
 +(v-u-w)^3+(w-u-v)^3\bigr].
 \label{eq:th-cubic-polarization}
\end{align}
Estas identidades algebraicas permiten exhibir canales escalares de cuadrados
o cubos cuando ello simplifica
\(W(z)\) y los coeficientes de Fourier.

La minimalidad no se decide contando monomios solamente. Si un término
cuadrático escalar tiene la forma \(\psi(r^{\mathsf T}x)\) con
\(\psi(s)=\kappa s^2\), su Hessiano es
\[
 \nabla^2[\kappa(r^{\mathsf T}x)^2]=2\kappa rr^{\mathsf T},
\]
de rango a lo sumo uno. Un Hessiano de rango mayor que uno prueba que ese
término no puede reducirse a un solo canal cuadrático.

\subsection{Balance armónico sesgado y multiharmónico}
\label{sec:multiharmonic}

Se aproxima el estado estacionario por
\begin{equation}
 x_H(t)=X_0+
 \sum_{h=1}^{H}\left(X_he^{\mathrm i h\omega t}
 +\overline{X_h}e^{-\mathrm i h\omega t}\right).
 \label{eq:th-fourier-state}
\end{equation}
Sea \(y_H=Cx_H\). Los coeficientes exactos del diccionario son
\begin{equation}
 \widehat\Phi_h(X_0,\ldots,X_H)
 =\frac1{2\pi}\int_0^{2\pi}
 \Phi\!\bigl(y_H(\theta)\bigr)e^{-\mathrm i h\theta}\,d\theta.
 \label{eq:th-phi-fourier}
\end{equation}
La reconstrucción de estos coeficientes es explícita para diccionarios
polinómicos. Si \(Y_k=CX_k\), con \(X_{-k}=\overline{X_k}\) y
\(X_k=0\) para \(|k|>H\), el producto de dos componentes satisface
\begin{equation}
 \widehat{(y_i y_j)}_h
 =\sum_{k=-H}^{H}(Y_k)_i(Y_{h-k})_j.
 \label{eq:th-quadratic-convolution}
\end{equation}
La identidad se obtiene al multiplicar las dos sumas de Fourier y agrupar
los términos con exponente \(e^{\mathrm ih\theta}\). Para un producto
cúbico, el mismo argumento da
\begin{equation}
 \widehat{(y_i y_j y_\ell)}_h
 =\sum_{k=-H}^{H}\sum_{m=-H}^{H}
 (Y_k)_i(Y_m)_j(Y_{h-k-m})_\ell.
 \label{eq:th-cubic-convolution}
\end{equation}
En particular, el coeficiente medio de un producto cuadrático contiene
\((Y_0)_i(Y_0)_j+2\operatorname{Re}\sum_{k=1}^{H}
(Y_k)_i\overline{(Y_k)_j}\). Los productos generan así DC y armónicos
hasta \(2H\) o \(3H\), respectivamente; truncar de nuevo a \(H\) deja
los residuos de \cref{eq:th-physical-harmonic-tail}.

Al igualar coeficientes en \cref{eq:th-mimo}, se obtiene
\begin{align}
 0&=AX_0+d+B\widehat\Phi_0,
 \label{eq:th-dc-balance}\\
 \bigl[(\mathrm ih\omega)^qI-A\bigr]X_h
 &=B\widehat\Phi_h,\qquad h=1,\ldots,H.
 \label{eq:th-harmonic-balance}
\end{align}
El término constante \(d\) sólo contribuye a \(h=0\); la derivada
estacionaria de esa componente es cero. Para \(h\ne0\), el operador de
Liouville--Weyl multiplica el coeficiente por \((\mathrm ih\omega)^q\),
mientras \(A\) actúa componente a componente. Estas tres contribuciones
producen las ecuaciones anteriores. La rama principal es
\[
 (\mathrm ih\omega)^q=(h\omega)^q
 \left[\cos\frac{q\pi}{2}+\mathrm i\sin\frac{q\pi}{2}\right].
\]
El sistema real se obtiene separando partes reales e imaginarias y fijando una
fase, por ejemplo \(\Im(e_j^{\mathsf T}X_1)=0\). Sin esa condición, la
invariancia temporal deja una familia continua de soluciones equivalentes.

Una no linealidad cuadrática excitada por un fundamental genera DC y segundo
armónico; una cúbica genera fundamental y tercero. Por eso \(H=1\) no cierra
exactamente Lorenz, RF, Genesio--Tesi o el jerk cuadrático. Los mínimos
estructurales \(H\ge2\) y \(H\ge3\) son puntos de partida, no garantías de
convergencia. Para \(|h|>H\), el coeficiente omitido del residual de la
ecuación es \(R_h=-B\widehat\Phi_h\). Su medida física es
\begin{equation}
 \epsilon_H^2=\sum_{|h|>H}\|B\widehat\Phi_h\|_2^2,
 \qquad
 \eta_H^2=\sum_{|h|>H}\|\widehat\Phi_h\|_2^2.
 \label{eq:th-physical-harmonic-tail}
\end{equation}
La cantidad \(\eta_H\) sólo mide la cola del diccionario. Si la composición
\(\Phi(Cx_H)\) es cuadrado integrable, Parseval identifica
\(\epsilon_H\) con la norma cuadrática media de la parte omitida del
residual físico. Se tiene
\(\epsilon_H\le\|B\|_{2\to2}\eta_H\), pero la recíproca puede fallar por
cancelación entre canales.

Por ejemplo, para \(q=1\),
\[
 A=\begin{pmatrix}0&-1\\1&0\end{pmatrix},\qquad
 B=\begin{pmatrix}1&-1\\0&0\end{pmatrix},\qquad
 \Phi(x)=(x_1^2,x_1^2)^{\mathsf T},
\]
la contribución \(B\Phi\) se anula idénticamente y
\(x(t)=(\cos t,\sin t)^{\mathsf T}\) es una solución exacta con \(H=1\).
No obstante, \(\eta_1^2=1/4\), mientras \(\epsilon_1=0\). Por ello una
cota impuesta sólo a \(\eta_H\) puede rechazar una solución exacta y depende
del escalamiento o redundancia del diccionario.

Para aceptar el cierre se controlan, en escalas físicas declaradas, el
residual DC \(R_0=-AX_0-d-B\widehat\Phi_0\), los residuos retenidos
\(R_h=[(\mathrm ih\omega)^qI-A]X_h-B\widehat\Phi_h\), la fase y
\(\epsilon_H\), además del error de cuadratura. Al aumentar \(H\) se
comprueba la estabilización con tolerancias predeclaradas. Si la no linealidad
genera infinitos armónicos, una FFT finita necesita además una cota de cola y
control de aliasing. Para \(0<q<1\), estos residuos pertenecen al cierre
estacionario de Liouville--Weyl; su pequeñez no demuestra una solución
periódica exacta del PVI causal de Caputo.

\begin{proposalbox}
\textbf{Admisión de semillas.} Acoplar el certificado D0/T1/M1 con
\cref{eq:th-dc-balance,eq:th-harmonic-balance}; una semilla M1 se admite como
candidata cuando satisface fase, cierre DC, residual físico \(\epsilon_H\) de
armónicos omitidos y
residual del campo original.
\end{proposalbox}

\subsection{Homotopía causal de Caputo con una sola historia}
\label{sec:single-history}

Para \(0<q<1\), una continuación parametrizada por
\(\varepsilon(t)\) debe interpretarse como el único problema no autónomo
\begin{equation}
 x(t)=x_0+\frac1{\Gamma(q)}
 \int_{t_0}^{t}(t-\tau)^{q-1}
 F_{\varepsilon(\tau)}(x(\tau))\,d\tau.
 \label{eq:th-single-history}
\end{equation}
Si \(\varepsilon\) cambia en tiempos \(t_j\), las contribuciones de
\([t_0,t_j]\) permanecen en el integral para todo \(t>t_j\). Reiniciar en cada
etapa con
\[
 x(t_j)=x_j,\qquad
 {}^{\mathrm C}D_{t_j}^q x=F_{\varepsilon_j}(x)
\]
elimina esa historia y define otro PVI.

\begin{proposition}[distinción de tres objetos]
\label{prop:history-three-objects}
En una continuación Caputo deben distinguirse:
\begin{enumerate}
 \item la trayectoria no autónoma completa de
 \cref{eq:th-single-history};
 \item el segmento posterior a \(t_J\) con
 \(\varepsilon(t)=\varepsilon_\star\), pero con la historia heredada desde
 \(t_0\);
 \item la solución autónoma reinicializada en \(t_J\), con terminal inferior
 \(t_J\) y estado inicial \(x(t_J)\).
\end{enumerate}
Los objetos 2 y 3 tienen el mismo campo instantáneo y el mismo estado en
\(t_J\), pero en general no satisfacen la misma ecuación integral.
\end{proposition}

\begin{proof}
Para el objeto 2, separar el integral en \(t_J\) produce
\[
 \int_{t_0}^{t_J}(t-\tau)^{q-1}
 F_{\varepsilon(\tau)}(x(\tau))\,d\tau
 +\int_{t_J}^{t}(t-\tau)^{q-1}
 F_{\varepsilon_\star}(x(\tau))\,d\tau.
\]
El primer término depende de \(t\) y no es una constante absorbible en la
condición inicial. El objeto 3 contiene únicamente el segundo integral con
otra trayectoria. Salvo \(q=1\), una historia nula o una coincidencia
especial, las soluciones difieren.
\end{proof}

Durante el transporte se conserva el objeto 2 y se
evalúa además el objeto 3 como prueba independiente del campo final. Sólo el
objeto 3 responde directamente a «¿qué hace el PVI autónomo final desde ese
estado sin memoria heredada?». La comparación debe conservar ambos resultados.

\begin{proposalbox}
\textbf{Hipótesis de continuación.} El transporte causal conserva la historia
y mantiene fijo el campo en la etapa final. La comparación de reinicios usa
el mismo terminal inferior. La dependencia continua respecto de
\(\varepsilon\) requiere hipótesis de existencia, unicidad y Lipschitz
uniformes para la ecuación de Volterra.
\end{proposalbox}

\ifHAFOBasinMaterial
\subsection{Cero contactos, primer contacto y alcance estadístico}
\label{sec:zero-contact}

Si las condiciones se extraen de manera i.i.d. de una distribución declarada
y cada ensayo tiene probabilidad de contacto \(p\), observar cero contactos en
\(N\) ensayos da el límite superior unilateral exacto
\begin{equation}
 p\leq1-\alpha^{1/N}
 \label{eq:th-binomial-upper}
\end{equation}
con confianza \(1-\alpha\). Para \(\alpha=0.05\), el límite es aproximadamente
\(3/N\) cuando \(N\) es grande. Esta interpretación no se aplica
automáticamente a mallas deterministas ni a trayectorias dependientes. En esos
casos el conteo es descriptivo y la afirmación correcta conserva la geometría
de muestreo.

El protocolo recomendado informa para cada equilibrio y radio:
\[
 (N_{\rm contacto},N_{\rm otro},N_{\rm divergente},
 N_{\rm ambiguo},N_{\rm fallo})
\]
y busca el primer radio con contacto. El primer contacto refuta la separación
para ese punto; cero contactos no prueba que ninguna vecindad intercepte la
cuenca.
\fi

\subsection{Taxonomía de rutas y protocolo L0--L6}
\label{sec:l0-l6}

La selección estructural precede a la integración:

\begin{description}[style=nextline,leftmargin=2.7em,labelwidth=2.2em]
 \item[D0] realización Lur'e escalar exacta y cierre directo;
 \item[T1] transformación afín u observable que produce D0;
 \item[A1] ruta alternativa activada por una obstrucción demostrada
 (mapa de conmutación, ganancia cero, frontera de Hopf);
 \item[M1] realización multicanal y balance sesgado/multiharmónico.
\end{description}

Los niveles de evaluación dinámica son:
\begin{description}[style=nextline,leftmargin=2.7em,labelwidth=2.2em]
 \item[L0] fuente, ecuaciones y parámetros fijados;
 \item[L1] álgebra y equilibrios verificados;
 \item[L2] semilla con residual y origen registrado;
 \item[L3] continuación con trazas y memoria declarada;
 \item[L4] recurrencia y clasificación dinámica;
 \item[L5] sondeos de todos los equilibrios y primer contacto;
 \item[L6] reproducción independiente y robustez de paso/horizonte.
\end{description}

La taxonomía no demuestra existencia de un atractor. Su función es impedir que
una semilla, un espectro o una figura aislada se presente como validación
completa.
